% Options for packages loaded elsewhere
\PassOptionsToPackage{unicode}{hyperref}
\PassOptionsToPackage{hyphens}{url}
\PassOptionsToPackage{dvipsnames,svgnames,x11names}{xcolor}
%
\documentclass[
  10pt,
]{article}
\usepackage{amsmath,amssymb}
\usepackage{iftex}
\ifPDFTeX
  \usepackage[T1]{fontenc}
  \usepackage[utf8]{inputenc}
  \usepackage{textcomp} % provide euro and other symbols
\else % if luatex or xetex
  \usepackage{unicode-math} % this also loads fontspec
  \defaultfontfeatures{Scale=MatchLowercase}
  \defaultfontfeatures[\rmfamily]{Ligatures=TeX,Scale=1}
\fi
\usepackage{lmodern}
\ifPDFTeX\else
  % xetex/luatex font selection
  \setmainfont[]{DejaVu Serif}
  \setsansfont[]{DejaVu Sans}
  \setmonofont[Scale=0.82]{DejaVu Sans Mono}
\fi
% Use upquote if available, for straight quotes in verbatim environments
\IfFileExists{upquote.sty}{\usepackage{upquote}}{}
\IfFileExists{microtype.sty}{% use microtype if available
  \usepackage[]{microtype}
  \UseMicrotypeSet[protrusion]{basicmath} % disable protrusion for tt fonts
}{}
\makeatletter
\@ifundefined{KOMAClassName}{% if non-KOMA class
  \IfFileExists{parskip.sty}{%
    \usepackage{parskip}
  }{% else
    \setlength{\parindent}{0pt}
    \setlength{\parskip}{6pt plus 2pt minus 1pt}}
}{% if KOMA class
  \KOMAoptions{parskip=half}}
\makeatother
\usepackage{xcolor}
\usepackage[margin=2.6cm]{geometry}
\usepackage{longtable,booktabs,array}
\usepackage{calc} % for calculating minipage widths
% Correct order of tables after \paragraph or \subparagraph
\usepackage{etoolbox}
\makeatletter
\patchcmd\longtable{\par}{\if@noskipsec\mbox{}\fi\par}{}{}
\makeatother
% Allow footnotes in longtable head/foot
\IfFileExists{footnotehyper.sty}{\usepackage{footnotehyper}}{\usepackage{footnote}}
\makesavenoteenv{longtable}
\usepackage{graphicx}
\makeatletter
\def\maxwidth{\ifdim\Gin@nat@width>\linewidth\linewidth\else\Gin@nat@width\fi}
\def\maxheight{\ifdim\Gin@nat@height>\textheight\textheight\else\Gin@nat@height\fi}
\makeatother
% Scale images if necessary, so that they will not overflow the page
% margins by default, and it is still possible to overwrite the defaults
% using explicit options in \includegraphics[width, height, ...]{}
\setkeys{Gin}{width=\maxwidth,height=\maxheight,keepaspectratio}
% Set default figure placement to htbp
\makeatletter
\def\fps@figure{htbp}
\makeatother
\setlength{\emergencystretch}{3em} % prevent overfull lines
\providecommand{\tightlist}{%
  \setlength{\itemsep}{0pt}\setlength{\parskip}{0pt}}
\setcounter{secnumdepth}{-\maxdimen} % remove section numbering
\ifLuaTeX
  \usepackage{selnolig}  % disable illegal ligatures
\fi
\IfFileExists{bookmark.sty}{\usepackage{bookmark}}{\usepackage{hyperref}}
\IfFileExists{xurl.sty}{\usepackage{xurl}}{} % add URL line breaks if available
\urlstyle{same}
\hypersetup{
  colorlinks=true,
  linkcolor={NavyBlue},
  filecolor={Maroon},
  citecolor={Blue},
  urlcolor={NavyBlue},
  pdfcreator={LaTeX via pandoc}}

\author{}
\date{}

\begin{document}

\hypertarget{monotone-growth-of-content-addressed-rule-closures}{%
\section{Monotone Growth of Content-Addressed Rule
Closures}\label{monotone-growth-of-content-addressed-rule-closures}}

\hypertarget{conservativity-history-independence-and-a-parsimony-trilemma-for-incrementally-generated-knowledge-structures}{%
\subsubsection{Conservativity, history independence, and a parsimony
trilemma for incrementally generated knowledge
structures}\label{conservativity-history-independence-and-a-parsimony-trilemma-for-incrementally-generated-knowledge-structures}}

\textbf{Behnam Sour, MD} · ORCID 0009-0002-3176-1956 Orthopaedic
Surgery, Shafa Yahyaeian Hospital, Tehran · Bone \& Joint Reconstruction
Research Center, Iran University of Medical Sciences

\begin{quote}
\textbf{Working paper · v2 · not peer-reviewed.} Formal framework with a
three-rung verification ladder: a machine-checked reference
implementation (\texttt{growth\_check.py}; 15/15 ledger checks), a
randomized property suite (\texttt{test\_growth\_properties.py}; 14/14
properties on \textasciitilde840 adversarial instances with shrinking),
and a self-contained Lean 4 formalization (\texttt{GrowthClosure.lean};
kernel-checked, \texttt{sorry}-free, no \texttt{Classical.choice}).
Every numeric value, address, fingerprint, benchmark time, and corpus
statistic in this document is produced by the accompanying scripts, not
illustrated by hand. Relative to v1, this revision repositions Theorem
11 against the CALM theorem following a prior-art sweep (§5, §9), adds
empirical scale measurements and the verification ladder (§8), and
replaces informal attributions with an existence-verified bibliography.
\end{quote}

\begin{center}\rule{0.5\linewidth}{0.5pt}\end{center}

\hypertarget{abstract}{%
\subsection{Abstract}\label{abstract}}

We study the \emph{growth problem} for content-addressed rule closures:
a knowledge structure is generated as the least fixpoint \(K(S)\) of a
finite family of typed generating rules over a seed set \(S\) of atoms,
every part is identified by a recursive cryptographic hash of its
derivation content, and the question is what happens when the seed
grows. Under four standing hypotheses --- finitary rules,
payload-determined side conditions, stratified acyclic generation, and
content addressing --- we prove that growth is \textbf{conservative}
(\(S \subseteq S'\) implies \(K(S)\) embeds in \(K(S')\) at identical
addresses, so growth is auditable as a Merkle address-set difference),
\textbf{incremental} (\(K(S \cup D) = K(K(S) \cup D)\), computable by
frontier-local, semi-naive iteration with no deletion or rederivation),
\textbf{history-independent} (the canonical fingerprint depends only on
the union of ingested batches, verified over all 13 ordered set
partitions of a 3-element delta), and \textbf{continuous}
(\(K(\bigcup_i S_i) = \bigcup_i K(S_i)\) for directed families,
licensing streaming semantics). We then prove the qualitative dynamics:
admissibility is stable but inertness is not, so seed growth is the
unique, append-only repair channel for ill-formed structures
(absorption); well-formedness itself is \emph{not} monotone and must be
re-checked by one breadth-first search per ingestion; and the first
Betti number of the underlying graph is monotone non-decreasing along
generation, with per-instance increment exactly \(m - j\). Finally we
isolate the cost of these guarantees: descriptor cohorts of size \(n\)
generate \(\sum_{m \ge m_0}\binom{n}{m}\) admissible bindings
(exponential; \(42\) at \(n=6\)), and we localize a known monotonicity
boundary --- the CALM theorem of Hellerstein and of Ameloot, Neven \&
Van den Bussche --- in this rule format: no payload-determined side
condition can store exactly the \(\subseteq\)-maximal bindings at every
seed (Theorem 11, with a kernel-checked decidable witness), giving a
trilemma: conservativity, history independence, and extensional maximal
parsimony are jointly unsatisfiable. Three resolutions are developed
(subsumption-quotient views, versioned cohort objects, and two-phase
closure-then-selection), all instances of one design principle:
\emph{keep identity and storage monotone and complete; spend parsimony
in views.} The closure algebra and the impossibility witness are
machine-proved in a self-contained, constructive Lean 4 development; on
a natural-language corpus the clique blow-up is real (one 268-atom
cohort alone implies \(\sim\!10^{80.7}\) admissible bindings), and a
fixed 20-atom delta integrates into bases up to 8,000 atoms
\(6\)--\(10\times\) faster incrementally than by batch reclosure.

\begin{center}\rule{0.5\linewidth}{0.5pt}\end{center}

\hypertarget{introduction}{%
\subsection{1. Introduction}\label{introduction}}

A generated, content-addressed knowledge structure faces a constraint
that hand-built structures do not: it must integrate new material
\textbf{without rebuilding and without rewriting history}. If every
part's identity is a hash of its derivation content, then any update
discipline that retracts or rewrites parts breaks every reference ever
issued to them; and if the structure is the value of a fixpoint, any
update discipline that recomputes from scratch wastes everything already
derived. The growth problem is to characterize exactly which guarantees
an incrementally fed closure can have, at what cost, and where the
unavoidable trade-off sits.

This paper answers the question for a general class of systems: typed
directed graphs generated as least fixpoints of finitary rule signatures
with payload-determined side conditions, under recursive content
addressing. The class includes, as instances, Datalog-style saturation
with content-addressed Skolemization, Merkle-DAG knowledge stores fed by
aggregation rules, and compression-gated (MDL-admissible) structure
induction over growing corpora. The theorems do not depend on any
particular type tower; a small four-level demonstration tower is used
only to make every claim machine-checkable.

\textbf{Contributions.}

\begin{enumerate}
\def\labelenumi{\arabic{enumi}.}
\tightlist
\item
  \textbf{Growth theorems} (§3). Conservativity with address-stable
  embedding (Thm. 2); the incremental-closure identity
  \(K(S\cup D)=K(K(S)\cup D)\) (Thm. 3); frontier locality,
  i.e.~correctness of semi-naive delta evaluation with no deletion phase
  (Thm. 4); ingestion confluence --- the fingerprint is a function of
  corpus content, not ingestion history (Thm. 5); and Scott-continuity
  in the seed, licensing streaming (Thm. 6).
\item
  \textbf{Qualitative dynamics} (§4). Stability of activation
  vs.~instability of inertness; absorption as the unique append-only
  repair mechanism; non-monotonicity of well-formedness with a one-BFS
  re-check discipline (Prop. 8); monotone growth of the first Betti
  number with the exact per-instance law \(\Delta\beta_1 = m - j\)
  (Prop. 9).
\item
  \textbf{A localized impossibility boundary} (§5). The
  descriptor-clique blow-up is quantified (Prop. 10); Theorem 11 proves
  that no payload-determined side condition realizes exactly the
  maximal-binding field at every seed --- positioned, following the
  prior-art sweep, as a content-addressed, MDL-flavored instance of the
  CALM boundary (Hellerstein 2010; Ameloot, Neven \& Van den Bussche
  2013; Hellerstein \& Alvaro 2020) rather than a new phenomenon.
  Corollary 12 states the resulting trilemma, whose conservativity
  vertex --- address-stable embedding --- has no CALM counterpart, with
  machine-checked witnesses including a real retracted address.
\item
  \textbf{Resolutions} (§6). Three architectures that each sacrifice a
  different vertex of the trilemma, with their guarantees stated
  precisely.
\item
  \textbf{A three-rung verification ladder and empirical scale}
  (§7--§8). A pure-stdlib reference implementation reproduces every
  quantitative claim (real SHA-256 fingerprints, confluence over all 13
  ingestion histories, the retraction counterexample, the 42-block
  clique, \(\beta_1 = 150\), an \(11\times11\) threshold grid); a
  randomized property suite re-tests the universal statements on
  \textasciitilde840 adversarial instances; a self-contained Lean 4
  development kernel-checks the closure algebra (constructively) and the
  impossibility witness (by \texttt{decide}); and two measurements
  quantify delta-integration cost and the clique blow-up on a natural
  corpus.
\end{enumerate}

\textbf{Relation to prior frameworks.} The setting axiomatizes, in
neutral terms, the growth behaviour of any typed generative-closure
system satisfying H1--H4 below. The components are individually
classical --- least fixpoints, semi-naive evaluation, the chase, MDL,
Merkle stores; §9 gives the lineage. The contribution claimed is the
growth analysis itself: the four guarantees, the impossibility result
that bounds them, and the verified implementation.

\begin{center}\rule{0.5\linewidth}{0.5pt}\end{center}

\hypertarget{formal-setting}{%
\subsection{2. Formal setting}\label{formal-setting}}

\hypertarget{atoms-types-parts}{%
\subsubsection{2.1 Atoms, types, parts}\label{atoms-types-parts}}

\textbf{Definition 1 (atom).} An \emph{atom} is a pair \(a = (p, h)\)
where \(p\) is a payload string and \(h = \mathrm{SHA\text{-}256}(p)\)
is its address. Each atom carries a \emph{descriptor} \(d(a)\), a
designated payload field naming the shared feature class it
instantiates. Atoms are the generators: no rule outputs type
\(\mathsf{Atom}\).

\textbf{Definition 2 (type poset).} \(T\) is a finite graded poset with
grade \(\rho : T \to \mathbb{N}\), \(\mathsf{Atom}\) the unique element
of grade \(0\), and a terminal constant \(\top\) above all types,
reached only by governance edges. The \emph{demonstration tower} used
throughout for machine checks is

\[\mathsf{Atom}\,(0) \;<\; \mathsf{Block}\,(1) \;<\; \mathsf{Section}\,(2) \;<\; \mathsf{Root}\,(3); \qquad \top \text{ constant}.\]

\textbf{Definition 3 (\(\Sigma\)-graph).} A \(\Sigma\)-graph is a typed
directed graph \(G = (V, E, \tau)\), \(\tau : V \to T\), in which every
non-atom vertex is the output of exactly one instance of a rule from a
finite signature \(\Sigma\), with edges from the instance's premises to
its output, plus governance edges into \(\top\).

\textbf{Definition 4 (rule).} A rule \(R \in \Sigma\) is a tuple
\(R = (\mathrm{in}_R, \mathrm{out}_R, \varphi_R, \pi_R)\): a finite list
of input types, an output type, a side condition \(\varphi_R\), and a
payload constructor \(\pi_R\). The demonstration signature:

\begin{longtable}[]{@{}
  >{\raggedright\arraybackslash}p{(\columnwidth - 6\tabcolsep) * \real{0.2500}}
  >{\raggedright\arraybackslash}p{(\columnwidth - 6\tabcolsep) * \real{0.2500}}
  >{\raggedright\arraybackslash}p{(\columnwidth - 6\tabcolsep) * \real{0.2500}}
  >{\raggedright\arraybackslash}p{(\columnwidth - 6\tabcolsep) * \real{0.2500}}@{}}
\toprule\noalign{}
\begin{minipage}[b]{\linewidth}\raggedright
rule
\end{minipage} & \begin{minipage}[b]{\linewidth}\raggedright
typing
\end{minipage} & \begin{minipage}[b]{\linewidth}\raggedright
side condition
\end{minipage} & \begin{minipage}[b]{\linewidth}\raggedright
kind
\end{minipage} \\
\midrule\noalign{}
\endhead
\bottomrule\noalign{}
\endlastfoot
\(\mathsf{bind}\) & \(\mathsf{Atom}^m \to \mathsf{Block}\) &
\(m \ge 3\), shared descriptor \(s\), \(\Delta(s,m) > 0\) & ascent \\
\(\mathsf{lift}\) & \(\mathsf{Block} \to \mathsf{Section}\) & --- &
ascent \\
\(\mathsf{frame}\) & \(\mathsf{Section} \to \mathsf{Root}\) & --- &
ascent \\
\(\mathsf{close}\) & \(\mathsf{Root} \to \top\) & --- & governance \\
\end{longtable}

Here \(\Delta(s,m) = (m-1)\,b(s) - o - m r\) is the two-part-code (MDL)
saving of storing the shared descriptor once (at codebook overhead
\(o\), descriptor information \(b(s)\) bits, reference cost \(r\) per
use); a binding is \emph{admissible} iff it strictly compresses,
\(\Delta > 0\). The default code parameters in the ledger are \(b = 2\),
\(o = 2\), \(r = 0\), giving MDL threshold
\(m_0 = \lfloor o/b \rfloor + 2 = 3\) (verified on an \(11\times 11\)
grid, §7H).

\hypertarget{standing-hypotheses}{%
\subsubsection{2.2 Standing hypotheses}\label{standing-hypotheses}}

All results assume:

\begin{itemize}
\tightlist
\item
  \textbf{H1 (finitarity).} Every rule instance has finitely many
  premises.
\item
  \textbf{H2 (payload determinism).} Every side condition \(\varphi_R\)
  --- including the descriptor function, code lengths, and \(\Delta\)
  --- is a function of the premise payloads alone, fixed independently
  of the ambient graph.
\item
  \textbf{H3 (stratified acyclic generation).} Ascent rules strictly
  increase grade; the generating fragment below any part is a finite
  DAG; \(\top\) is a constant identifier outside the hash recursion, so
  governance edges create no content cycle.
\item
  \textbf{H4 (content addressing).} Addresses extend recursively from
  atoms to all parts:
\end{itemize}

\[\mathrm{addr}(v) = \begin{cases} \mathrm{SHA\text{-}256}(p) & v = (p,h) \text{ an atom},\\[2pt] \mathrm{SHA\text{-}256}\big(\mathrm{kind}(v) \,\|\, \mathrm{rule}(v) \,\|\, \mathrm{descr}(v) \,\|\, \mathrm{sort}\{\mathrm{addr}(c) : c \in \mathrm{in}(v)\}\big) & \text{otherwise,} \end{cases}\]

with \(\mathrm{addr}(\top) = \mathrm{SHA\text{-}256}(\texttt{"TOP"})\)
fixed. Parts are identified by address: duplicate derivation content
collapses to one node.

H2 is the load-bearing hypothesis. In a fully context-sensitive MDL code
the cost of a part depends on the global codebook, so admissibility can
flip as the field grows; restricting to payload-determined codes is
precisely what purchases monotonicity (Lemma A below) --- and §5 proves
that this purchase is unavoidable.

\hypertarget{closure-as-a-least-fixpoint}{%
\subsubsection{2.3 Closure as a least
fixpoint}\label{closure-as-a-least-fixpoint}}

Fix a countable atom universe; since \(\Sigma\) is finite and rules are
finitary, the universe \(\mathcal{U}\) of derivable records is
countable, and \((\mathcal{P}(\mathcal{U}), \subseteq)\) is a complete
lattice.

\textbf{Definition 5 (one-step expansion).} \(\Phi(X)\) is the set of
outputs of all rule instances whose premises lie in \(X\) and whose side
conditions hold. The \emph{inflationary} form is
\(\hat\Phi(X) = X \cup \Phi(X)\).

\textbf{Lemma A (persistence).} Under H2, if a side condition holds for
premises lying in \(X\) and \(X \subseteq Y\), it holds for the same
premises in \(Y\). \emph{Proof.} The condition is a predicate on premise
payloads; payloads do not change when the ambient set grows, and the
premises persist into \(Y\). \(\blacksquare\)

\textbf{Lemma B (monotonicity, continuity).} \(\Phi\) and \(\hat\Phi\)
are monotone; under H1 they are Scott-continuous:
\(\hat\Phi(\bigcup_i X_i) = \bigcup_i \hat\Phi(X_i)\) for directed
families. \emph{Proof.} Monotonicity from Lemma A. Continuity: any
output of \(\hat\Phi(\bigcup_i X_i)\) has finitely many premises, all in
a single \(X_j\) by directedness. \(\blacksquare\)

\textbf{Lemma C (existence, constructivity).} For any seed \(S\) the
closure \(K(S) = \mu X.\,(S \cup \Phi(X))\) exists and equals
\(\bigcup_{k \ge 0} \hat\Phi^k(S)\). \emph{Proof.} Knaster--Tarski gives
the least fixpoint of the monotone \(X \mapsto S \cup \Phi(X)\) on the
complete lattice; Kleene's theorem with Lemma B gives the constructive
\(\omega\)-union. \(\blacksquare\)

\hypertarget{canonical-form-and-fingerprint}{%
\subsubsection{2.4 Canonical form and
fingerprint}\label{canonical-form-and-fingerprint}}

\textbf{Definition 6 (canonical serialization).} \(\sigma(G)\) lists
every part as the record
\(\langle \mathrm{addr}, \mathrm{kind}, \mathrm{rule}, \mathrm{descr}, \text{sorted child addresses} \rangle\),
sorted by (grade, kind, address); governance edges are appended in
sorted order. The \emph{fingerprint} is
\(\mathrm{fp}(G) = \mathrm{SHA\text{-}256}(\sigma(G))\).

\textbf{Lemma D (order independence).} \(\mathrm{addr}(v)\) depends only
on the derivation content of \(v\), not on when \(v\) was derived; hence
the record set of \(K(S)\), and therefore \(\sigma\) and
\(\mathrm{fp}\), are independent of the order in which the fixpoint
iteration fires rules. \(\blacksquare\)

The fingerprint of the worked five-atom field of §7A is

\begin{verbatim}
1066bddea3ff771058f4af6afc828efbeb29f2968cf3c68cce02110cf64d9662
\end{verbatim}

with all part addresses pairwise distinct (injectivity checked,
including \(\top\)).

\begin{center}\rule{0.5\linewidth}{0.5pt}\end{center}

\hypertarget{growth-theorems}{%
\subsection{3. Growth theorems}\label{growth-theorems}}

Throughout, \(S \subseteq S'\) are seeds of atoms and \(D\) is a finite
\emph{delta} of new atoms.

\textbf{Theorem 1 (closure operator).} \(K\) is a closure operator on
seeds and on \(\Sigma\)-graphs: extensive (\(X \subseteq K(X)\)),
monotone (\(X \subseteq Y \Rightarrow K(X) \subseteq K(Y)\)), idempotent
(\(K(K(X)) = K(X)\)). \emph{Proof.} Extensivity by definition of the
fixpoint. Monotonicity: larger seeds give larger values of the monotone
map at every Kleene stage (Lemma B). Idempotence: \(K(X)\) is
\(\Sigma\)-closed, \(\Phi(K(X)) \subseteq K(X)\), so re-closing adds
nothing. \(\blacksquare\)

\textbf{Theorem 2 (conservativity; growth is a Merkle diff).} If
\(S \subseteq S'\) then: (i) \(K(S) \subseteq K(S')\) --- no part is
ever retracted; (ii) every part of \(K(S)\) occurs in \(K(S')\) at the
\emph{identical address}, so \(K(S)\) is an address-stable sub-DAG of
\(K(S')\); (iii) the record set of \(\sigma(K(S))\) is a subset of that
of \(\sigma(K(S'))\), and the growth is exactly the address-set
difference \(\mathrm{addr}(K(S')) \setminus \mathrm{addr}(K(S))\).
\emph{Proof.} (i) is Theorem 1's monotonicity, which rests on Lemma A:
under H2 a new atom can never invalidate an old instance's side
condition, and closure only adds. (ii) Addresses (H4) are functions of
derivation content alone, which is unchanged by ambient growth. (iii)
immediate from (i)--(ii) and Definition 6. \(\blacksquare\)

\emph{Consequences.} Old references never dangle;
checkpoint-to-checkpoint growth is auditable by comparing two sorted
address sets, with no replay; and Theorem 5 below makes the checkpoints
themselves canonical. (Descent is also stable: since no rule outputs
\(\mathsf{Atom}\), the atoms of \(K(S')\) are exactly \(S'\), so reading
the seed back remains exact under growth.)

\textbf{Theorem 3 (incremental closure).}
\(K(S \cup D) = K(K(S) \cup D)\). \emph{Proof.} (\(\supseteq\))
\(S \cup D \subseteq K(S) \cup D \subseteq K(S \cup D)\) by extensivity
and monotonicity; applying the monotone idempotent \(K\):
\(K(K(S) \cup D) \subseteq K(K(S \cup D)) = K(S \cup D)\).
(\(\subseteq\)) From \(S \cup D \subseteq K(S) \cup D\), monotonicity
gives \(K(S \cup D) \subseteq K(K(S) \cup D)\). \(\blacksquare\)

One never recomputes from scratch: seed the delta into the
already-closed field and re-close.

\textbf{Theorem 4 (frontier locality; semi-naive correctness).} Define
the frontier iteration from a closed field: \(F_0 = D \setminus K(S)\);
\(F_{k+1} =\) the outputs, not already present, of rule instances with
all premises in \(K(S) \cup \bigcup_{j \le k} F_j\) and \textbf{at least
one premise in} \(\bigcup_{j\le k} F_j\). Then
\(K(S) \cup \bigcup_k F_k = K(S \cup D)\). Consequently the delta
closure is computable by firing only instances that touch the frontier;
there is no deletion phase and no rederivation. \emph{Proof.}
(\(\subseteq\)) every frontier output is derivable from \(K(S)\cup D\),
hence in \(K(S\cup D)\) (Thm. 3). (\(\supseteq\)) Key step: if an
instance has \emph{all} premises in \(K(S)\) and its side condition
holds, its output is in \(\Phi(K(S)) \subseteq K(S)\) already ---
closedness plus H2. So any part of \(K(S\cup D)\setminus K(S)\) is
produced by an instance with at least one premise outside \(K(S)\).
Induct on Kleene stage: a stage-\((k{+}1)\) part not in \(K(S)\) has all
premises at stage \(\le k\), each either in \(K(S)\) or (by induction)
in some \(F_j\); at least one premise is of the latter kind, so the
instance fires in the frontier iteration. \(\blacksquare\)

\emph{Cost.} With premise indexing by (descriptor, kind), work is
proportional to the size of the delta closure times rule fan-in --- the
growth region only. The contrast with classical incremental view
maintenance is structural: deletion-capable maintenance (e.g.~DRed) must
over-delete and rederive because facts can lose support; conservativity
(Thm. 2) makes support loss impossible, so the entire deletion machinery
vanishes.

\textbf{Theorem 5 (ingestion confluence; history independence).} Let
\(D\) be a finite delta and \((D_1, \dots, D_k)\) \emph{any} ordered
partition of \(D\) into nonempty batches. Iterating incremental closure,

\[X_0 = K(S), \qquad X_j = K(X_{j-1} \cup D_j),\]

yields \(X_k = K(S \cup D)\), and hence the identical canonical form and
fingerprint, for every partition and every order. Checkpoint
fingerprints are therefore well-defined functions of corpus
\emph{content}, not ingestion \emph{history}. \emph{Proof.} Induction on
\(j\) with Theorem 3:
\(X_j = K\big(K(S \cup D_1 \cup \dots \cup D_{j-1}) \cup D_j\big) = K(S \cup D_1 \cup \dots \cup D_j)\).
The fingerprint depends only on the record set of the final closure
(Lemma D). \(\blacksquare\)

\emph{Machine check (§7B).} For a 3-element delta there are exactly
\(13\) ordered set partitions (the ordered Bell number \(a(3) = 13\)).
All \(13\) histories --- one batch, three sequential singletons in all
orders, and every mixed batching --- produce \textbf{one} fingerprint,
equal to the single-batch fingerprint:

\begin{verbatim}
5c8942ed83c9b398554ea4898f04f4eebe28ac0796b6d099b9e67e49fe76d38a
\end{verbatim}

\textbf{Theorem 6 (continuity in the seed; streaming semantics).} For
any directed family of seeds \(\{S_i\}\),
\(\;K\big(\bigcup_i S_i\big) = \bigcup_i K(S_i)\). \emph{Proof.} By
Lemma B each \(\hat\Phi^k\) is Scott-continuous, so
\(\hat\Phi^k(\bigcup_i S_i) = \bigcup_i \hat\Phi^k(S_i)\) by induction
on \(k\); then
\(K(\bigcup_i S_i) = \bigcup_k \bigcup_i \hat\Phi^k(S_i) = \bigcup_i K(S_i)\),
the unions commuting. \(\blacksquare\)

The field of a growing corpus is the colimit of the stage fields: a
streaming ingester computes, at every moment, a genuine sub-object of
the eventual field, and nothing computed is ever wasted. (Theorems 2--6
are exactly the algebra of a grow-only replicated set over records:
merge is union, which is associative, commutative, and idempotent ---
confluence is convergence.)

\begin{center}\rule{0.5\linewidth}{0.5pt}\end{center}

\hypertarget{qualitative-dynamics-of-growth}{%
\subsection{4. Qualitative dynamics of
growth}\label{qualitative-dynamics-of-growth}}

\textbf{Proposition 7 (activation is stable; inertness is not).} If a
rule instance is admissible over \(S\), it is admissible over every
\(S' \supseteq S\) and its output persists (Lemma A, Thm. 2). The
converse fails: an atom with no admissible instance over \(S\) (an
\emph{inert} or \emph{orphan} atom) may acquire one over
\(S' \supset S\), namely when growth supplies the missing cohort
members. Growth is therefore the \textbf{unique repair channel}: the
graph is never edited; the seed is grown. Repair is append-only.

\textbf{Definition (well-formedness).} Let \(\widehat G\) be the
undirected graph underlying \(G\) (content plus governance edges). \(G\)
is \emph{well-formed} iff \(\widehat G\) is connected to \(\top\) and
\(\mathrm{ecc}_{\widehat G}(\top) \le h\), where \(h\) is the ascent
height of the tower (\(h = 4\) for the demonstration tower: the extremal
path is
\(\mathsf{Atom}\)--\(\mathsf{Block}\)--\(\mathsf{Section}\)--\(\mathsf{Root}\)--\(\top\)).
Decidable by one breadth-first search from \(\top\).

\textbf{Proposition 8 (well-formedness is not monotone).} Growth neither
preserves nor reflects well-formedness: (a) \emph{repair by absorption}
--- an ill-formed field with an orphan atom becomes well-formed when the
orphan's descriptor cohort reaches the admissibility threshold and
binds; (b) \emph{breakage by novelty} --- a well-formed field becomes
ill-formed when a delta introduces an atom whose descriptor cohort is
below threshold. Hence well-formedness is a \textbf{validation
predicate, not a generation gate}: the closure always exists and all of
§3 applies regardless; the discipline is one BFS re-check after every
ingestion. \emph{Proof.} (a) and (b) are realized by explicit fields;
the ledger exhibits (a). \(\blacksquare\)

\emph{Machine check (§7C).} Stage 1: cohort \(\{c_1,c_2,c_3\}\) on
\(D_1\) plus one orphan on \(D_2\) --- \textbf{ill-formed} (\(1\)
orphan; \(\mathrm{ecc}(\top)=4\) on the reached component). Stage 2: a
2-atom delta on \(D_2\) completes the cohort; the orphan binds, lifts,
frames, closes --- \textbf{well-formed}, \(0\) orphans,
\(\mathrm{ecc}(\top) = 4\). Stage-1 addresses are a subset of stage-2
addresses (conservativity), and the incremental fingerprint equals the
batch fingerprint:

\begin{verbatim}
stage 1   2b2e50445f9e1bebfb9a22de28b087b1424585deae6c3fa659b97e3d56b79866   ill-formed
stage 2   97f6082afaba7abafa97eb69a42f73989c4018e31d55f5ecbd6a7f77c80dd497   well-formed
\end{verbatim}

\textbf{Proposition 9 (the first Betti number is monotone; per-instance
law).} Let \(\beta_1 = |E| - |V| + C\) on \(\widehat G\) (\(C\) = number
of connected components). Along any generation or growth trace: (i)
adding a fresh atom: \(\Delta\beta_1 = 0\); (ii) adding a composite with
\(m\) premises spanning \(j\) distinct existing components:
\(\Delta V = 1\), \(\Delta E = m\), \(\Delta C = 1 - j\), hence
\(\Delta\beta_1 = m - j \ge 0\) (since \(j \le m\)); (iii) adding a
governance edge between existing nodes: \(\Delta\beta_1 = 0\) if it
merges two components, \(1\) if its endpoints were already connected.
Therefore \(\beta_1\) is non-decreasing under generation and under seed
growth; since every connected graph deformation-retracts onto a wedge of
\(\beta_1\) circles, \(\pi_1(\widehat G)\) is free of rank \(\beta_1\)
throughout, and the homotopy type only ever \emph{accumulates} circles.
\emph{Proof.} Direct computation of the Euler-type quantity in each
case; \(j \le m\) because each premise lies in some component and there
are \(m\) premises. \(\blacksquare\)

\emph{Machine check (§7G).} On the \(n=6\) clique field: \(V = 133\),
\(E = 282\), \(C = 1\), \(\beta_1 = 150\); a union-find replay of the
entire generation reproduces \(\beta_1 = 150\) and confirms
\(\Delta\beta_1 = m - j\) at every ascent instance. (Sanity:
\(V = 6 + 42 + 42 + 42 + 1\);
\(E = \sum_{m=3}^{6} m\binom{6}{m} + 42 + 42 + 42 = 156 + 126\).)

\begin{center}\rule{0.5\linewidth}{0.5pt}\end{center}

\hypertarget{the-parsimony-dilemma}{%
\subsection{5. The parsimony dilemma}\label{the-parsimony-dilemma}}

The guarantees of §3 are purchased by H2, and H2 has a price. This
section quantifies the price and proves it cannot be negotiated away
within the rule format.

\textbf{Proposition 10 (descriptor-clique blow-up).}
\(\Delta(s,m) = (m-1)b - o - mr\) is strictly increasing in \(m\) iff
\(b > r\). Hence once a descriptor pays at the threshold
\(m_0 = \lfloor o/b \rfloor + 2\) (for \(r = 0\); verified on the
\(11\times 11\) grid with the characterization
\(m_0 = 3 \iff o/b \in [1,2)\)), \textbf{every} subset of the cohort of
size \(\ge \max(3, m_0)\) is admissible. A cohort of \(n\) atoms
therefore generates

\[N(n) \;=\; \sum_{m = \max(3, m_0)}^{n} \binom{n}{m} \;=\; 2^n - 1 - n - \binom{n}{2} \quad (m_0 \le 3),\]

binding nodes --- exponential in \(n\) --- and each binding lifts and
frames, so the redundancy propagates upward linearly per binding.
\emph{Machine check (§7F).} \(N(3),N(4),N(5),N(6) = 1, 5, 16, 42\),
exactly matching the closure (\(42\) blocks \(\to\) \(42\) sections
\(\to\) \(42\) roots at \(n = 6\)).

The obvious remedy is to bind \textbf{only \(\subseteq\)-maximal}
premise sets per descriptor. The next theorem shows this remedy is
incompatible with the entire §3 package.

\textbf{Theorem 11 (impossibility of payload-determined maximality).}
There is no side condition \(\varphi\) satisfying H2 such that, for
every seed \(S\), the closure under \(\varphi\) contains exactly the
\(\subseteq\)-maximal admissible bindings of each descriptor cohort in
\(S\). \emph{Proof.} Take four atoms \(a_1, \dots, a_4\) sharing one
descriptor, with \(m_0 = 3\). Let \(t = (a_1, a_2, a_3)\). Over
\(S_3 = \{a_1,a_2,a_3\}\), \(t\) is the unique maximal binding, so
exactness forces \(\varphi(t) = \mathrm{true}\). Over
\(S_4 = S_3 \cup \{a_4\}\), the premises of \(t\) have the same
payloads, so by H2 \(\varphi(t) = \mathrm{true}\) still --- yet \(t\) is
no longer maximal in \(S_4\), so exactness forces \(t\) \emph{out}.
Contradiction. (If instead \(\varphi(t) = \mathrm{false}\), exactness
already fails at \(S_3\).) Maximality of a premise set is irreducibly a
property of the \emph{ambient seed}, hence inexpressible
payload-determined. \(\blacksquare\)

\textbf{Corollary 12 (the growth trilemma).} No system in this rule
format simultaneously satisfies:

\begin{itemize}
\tightlist
\item
  \textbf{(C) conservativity} --- stages embed address-stably, nothing
  is retracted (Thm. 2);
\item
  \textbf{(H) history independence} --- the stored state is a function
  of seed content only (Thm. 5);
\item
  \textbf{(P) extensional maximal parsimony} --- the stored state
  contains exactly the maximal bindings.
\end{itemize}

Indeed (C)∧(P) is already contradictory: by Theorem 11's witness, the
maximal-binding states at \(S_3 \subset S_4\) satisfy
\(M(S_3) \not\subseteq M(S_4)\) --- the binding \(B_{123}\) is
\emph{retracted} --- so any operator \(\Phi_{\max}\) storing exactly
\(M(\cdot)\) is non-monotone, loses the Knaster--Tarski/Kleene route
entirely, and dangles every reference into the retracted cone. The
complete closure achieves (C)∧(H); §6's quotient view achieves (H)∧(P).

\emph{Machine check (§7D).} \(\Phi_{\max}(S_3)\) contains \(B_{123}\);
\(\Phi_{\max}(S_4)\) does not. The retracted address is real:

\begin{verbatim}
retracted   362db034e4f12612265ac2e3...        (B_123, descriptor D1)
\end{verbatim}

and the \(\mathsf{Section}\) lifted from \(B_{123}\) at stage 1 now
references a part absent from the stage-2 state. On the identical
growth, the complete closure contains \(B_{123}\) at \textbf{both}
stages and is conservative.

\textbf{Positioning.} Theorem 11 should be read as a content-addressed,
MDL-flavored instance of a known boundary rather than a new phenomenon.
That extremal aggregation is non-monotone is classical: semantics for
aggregates inside recursion require monotonicity with respect to a
lattice order precisely because operators like ``maximal set'' break it
(Ross \& Sagiv 1992/1997). The general form of the boundary is the CALM
theorem (conjectured by Hellerstein 2010; proved by Ameloot, Neven \&
Van den Bussche 2013; surveyed in Hellerstein \& Alvaro 2020): a problem
has a consistent, coordination-free distributed implementation iff it is
monotone --- where ``consistency'' is confluence over input orderings,
our history independence, and where the coordination-free computations
are captured by \emph{oblivious} transducers, deciding without global
knowledge, exactly as our payload-determined side conditions decide from
premise payloads alone. Under this dictionary, Theorem 11 instantiates
the only-if direction for maximal bindings, and Corollary 12 adds the
element CALM does not speak of: identity. In a content-addressed store,
the cost of non-monotonicity is not abstract ``coordination'' but the
retraction of named, referenced parts. The contribution claimed here is
accordingly the trilemma formulation with address-stable conservativity,
the localization of the boundary at MDL codebook globality, and the
mechanized witness --- not the existence of the boundary itself.

\begin{center}\rule{0.5\linewidth}{0.5pt}\end{center}

\hypertarget{resolutions-where-to-spend-the-parsimony}{%
\subsection{6. Resolutions: where to spend the
parsimony}\label{resolutions-where-to-spend-the-parsimony}}

All three architectures below keep the complete monotone closure for
\emph{identity} and differ in what they \emph{store} or \emph{present}.
Each gives up exactly one vertex of the trilemma.

\hypertarget{subsumption-quotient-view-keeps-hp-spends-c-in-the-view-only}{%
\subsubsection{6.1 Subsumption-quotient view --- keeps (H)∧(P), spends
(C) in the view
only}\label{subsumption-quotient-view-keeps-hp-spends-c-in-the-view-only}}

Define on bindings: \(B \sqsubseteq B'\) iff
\(\mathrm{descr}(B) = \mathrm{descr}(B')\) and
\(\mathrm{premises}(B) \subseteq \mathrm{premises}(B')\). Strict
containment of finite sets is well-founded, so maximal elements exist.
\textbf{Store} the complete field \(K(S)\); \textbf{present} the
quotient view \(\pi(K(S))\): drop subsumed bindings and the parts
derived \emph{exclusively} from them (their derivation cones). Then:

\begin{itemize}
\tightlist
\item
  \(\pi\) is history-independent --- it is a function of \(K(S)\), which
  is (Thm. 5);
\item
  \(\pi\) is maximally parsimonious at the binding level;
\item
  \(\pi\) is \textbf{non-conservative across stages} --- and harmlessly
  so: retraction lives only in the view; every address persists in
  storage, so references issued against the field never dangle.
\end{itemize}

\emph{Machine check (§7E).} Cohort \(\{g_1,g_2,g_3\}\) at stage 1: field
\(1\) block, view \(1\) block (\(B_{123}\) visible). Add \(g_7\): field
\(5\) blocks, view \(1\) block (\(B_{1237}\)); \(B_{123}\) has left the
\textbf{view} but remains in the \textbf{field} at both stages.

\hypertarget{versioned-cohort-objects-keeps-p-and-a-chain-level-c-spends-node-level-c}{%
\subsubsection{6.2 Versioned cohort objects --- keeps (P) and a
chain-level (C), spends node-level
(C)}\label{versioned-cohort-objects-keeps-p-and-a-chain-level-c-spends-node-level-c}}

Store one object per descriptor cohort (payload = the sorted member
address list), with bindings materialized on demand; storage is linear
in atoms + descriptors, eliminating the \(2^n\) blow-up at rest. The
cohort payload changes with the seed, so node-level conservativity is
lost by construction; it is \textbf{recovered at the version level} by
an append-only, hash-linked chain of cohort versions: each new version
references its predecessor's address, the chain embeds under growth, and
only the \emph{head pointer} moves. Fingerprints must then name the
materialization policy, since the same chain admits multiple
materialized fields. This resolution is the storage-layer analogue of
the distribution-aware relaxations of CALM (Zinn, Green \& Ludäscher
2012; Ameloot, Ketsman, Neven \& Zinn 2015): granting the system
knowledge beyond individual payloads --- here, the cohort version chain
--- buys back behavior that obliviousness forbids.

\hypertarget{two-phase-monotone-closure-global-selection-keeps-ch-states-p-as-optimization}{%
\subsubsection{6.3 Two-phase: monotone closure + global selection ---
keeps (C)∧(H), states (P) as
optimization}\label{two-phase-monotone-closure-global-selection-keeps-ch-states-p-as-optimization}}

Phase 1: the complete monotone closure --- identity, storage, and all of
§3. Phase 2: global MDL selection of a minimum-cost covering subfamily
of bindings per checkpoint. Selecting the cheapest set of bindings whose
union covers the cohort is weighted set cover --- NP-hard, with the
classical greedy \(\ln n\)-approximation --- and is therefore run as a
\textbf{non-monotone view recomputed per checkpoint}, never as a storage
discipline. Conservativity of the candidate space (Thm. 2) means each
re-optimization warm-starts from the previous solution plus only the
frontier's new candidates.

\textbf{Design principle.} \emph{Keep identity and storage monotone and
complete; spend parsimony in views.} Monotone storage buys Theorems 2--6
--- stable references, semi-naive deltas, canonical checkpoints,
streaming. Parsimony is real but representational; the trilemma says it
cannot be made extensional in storage without forfeiting growth.

\begin{center}\rule{0.5\linewidth}{0.5pt}\end{center}

\hypertarget{worked-growth-run-the-machine-checked-ledger}{%
\subsection{7. Worked growth run --- the machine-checked
ledger}\label{worked-growth-run-the-machine-checked-ledger}}

Reference implementation: \texttt{growth\_check.py} (pure Python 3
standard library: \texttt{hashlib}, \texttt{itertools},
\texttt{collections}). Reproduce with
\texttt{python3\ growth\_check.py}. \textbf{All 15 checks pass.}
Demonstration tower as in §2; code parameters \(b = 2, o = 2, r = 0\),
so \(m_0 = 3\).

\textbf{A --- Primary worked field.} Seed: three atoms on \(D_1\), one
on \(D_2\), one on \(D_3\). Only the \(D_1\) cohort binds (size 3); the
closure carries it to \(\mathsf{Root}\) and closes to \(\top\). The
\(D_2\)/\(D_3\) atoms are inert orphans.

\begin{longtable}[]{@{}lllll@{}}
\toprule\noalign{}
part & kind & rule & descr & addr (first 12 hex) \\
\midrule\noalign{}
\endhead
\bottomrule\noalign{}
\endlastfoot
\(a\) & Atom & input & \(D_1\) & \texttt{0e695abf2881} \\
\(a\) & Atom & input & \(D_1\) & \texttt{3c7f6d70e099} \\
\(a\) & Atom & input & \(D_1\) & \texttt{c7b946f36c7d} \\
\(a\) & Atom & input & \(D_2\) & \texttt{f6f6adf60b22} \\
\(a\) & Atom & input & \(D_3\) & \texttt{6554a14f87c5} \\
\(B\) & Block & bind & \(D_1\) & \texttt{900b1087dd65} \\
\(S\) & Section & lift & \(D_1\) & \texttt{17cd9dd60db5} \\
\(R\) & Root & frame & \(D_1\) & \texttt{bf358e13a3c4} \\
\(\top\) & const & --- & --- & \texttt{07b4ed8e4e4e} \\
\end{longtable}

Fingerprint
\texttt{1066bddea3ff771058f4af6afc828efbeb29f2968cf3c68cce02110cf64d9662}.
Addresses injective (incl.~\(\top\)). \(\mathrm{ecc}(\top) = 4\) on the
reached component; \(2\) orphans \(\Rightarrow\) \textbf{ill-formed} ---
the picture Proposition 8 repairs.

\textbf{B --- Ingestion confluence.} Base \(\{b_1, b_2\}\) on \(D_1\)
(no binding yet); delta \(\{d_1, d_2, d_3\}\) on \(D_1\). All \(13\)
ordered set partitions of the delta, replayed by frontier-local
incremental closure, yield \textbf{1 distinct fingerprint}, equal to the
single-batch fingerprint \texttt{5c8942ed83c9...e76d38a}. The 5-atom
cohort closes to \(16 = \binom{5}{3} + \binom{5}{4} + \binom{5}{5}\)
blocks --- confluence and the clique semantics in one run.

\textbf{C --- Absorption.} Ill-formed (1 orphan) \(\to\) well-formed (0
orphans, \(\mathrm{ecc}(\top) = 4\)) via a 2-atom delta completing the
\(D_2\) cohort; stage-1 address set \(\subset\) stage-2 address set;
incremental \(=\) batch fingerprint. Hashes as in §4.

\textbf{D --- Non-monotone counterexample.}
\(\Phi_{\max}(S_3) = \{B_{123}\}\), \(\Phi_{\max}(S_4) = \{B_{1234}\}\):
\(B_{123}\) (addr \texttt{362db034e4f1…}) \textbf{retracted},
monotonicity check \texttt{False}; complete closure on the same growth:
\(B_{123}\) present at both stages, conservativity \texttt{True}.

\textbf{E --- Quotient view.} Field blocks \(1 \to 5\), view blocks
\(1 \to 1\) across the \(g_7\) ingestion; view non-conservative
\texttt{True}, field conservative \texttt{True}.

\textbf{F --- Clique counts.} \(n = 3,4,5,6 \Rightarrow 1, 5, 16, 42\)
blocks, each matching \(\sum_{m\ge3}\binom{n}{m}\); sections and roots
track blocks one-for-one.

\textbf{G --- Betti number.} \(n=6\) field:
\(V=133,\ E=282,\ C=1,\ \beta_1 = 150\); union-find replay matches;
\(\Delta\beta_1 = m - j\) at every ascent instance.

\textbf{H --- Threshold grid.} Over the \(11\times11\) integer grid
\((b, o) \in \{1..11\}^2\), \(r=0\): brute-force least admissible \(m\)
equals \(\lfloor o/b\rfloor + 2\) in every cell, and
\(m_0 = 3 \iff o/b \in [1, 2)\). Samples: \((o{=}2, b{=}2)\):
\(\Delta(2) = 0\), \(\Delta(3) = 2\), \(m_0 = 3\); \((o{=}4, b{=}2)\):
\(\Delta(3) = 0\), \(\Delta(4) = 2\), \(m_0 = 4\).

\textbf{Ledger summary.} A: injectivity \(\checkmark\) · B: 13 histories
→ 1 fingerprint \(\checkmark\), incremental = batch \(\checkmark\) · C:
absorption \(\checkmark\), conservativity \(\checkmark\), incremental =
batch \(\checkmark\) · D: \(\Phi_{\max}\) non-monotone \(\checkmark\),
complete closure conservative \(\checkmark\) · E: view non-conservative
\(\checkmark\), field conservative \(\checkmark\) · F: \(n{=}6 \to 42\)
\(\checkmark\) · G: replay = direct \(\checkmark\), per-instance law
\(\checkmark\) · H: grid \(\checkmark\), characterization \(\checkmark\)
--- \textbf{15/15 PASS.}

\begin{center}\rule{0.5\linewidth}{0.5pt}\end{center}

\hypertarget{empirical-scale-and-mechanized-verification}{%
\subsection{8. Empirical scale and mechanized
verification}\label{empirical-scale-and-mechanized-verification}}

All measurements in this section are produced by
\texttt{bench\_growth.py} and \texttt{corpus\_cohorts.py} (pure stdlib +
matplotlib) against the unmodified reference implementation; rerun both
to reproduce every number and figure.

\hypertarget{delta-integration-cost-theorem-4s-remark-measured}{%
\subsubsection{8.1 Delta integration cost --- Theorem 4's remark,
measured}\label{delta-integration-cost-theorem-4s-remark-measured}}

Random corpora with every descriptor cohort capped at 5 atoms (so the
complete closure stays small per cohort); a fixed delta of \(|D| = 20\)
atoms, half landing in existing cohorts to exercise absorption; median
of 5 runs. At every size the incremental and batch fingerprints are
asserted equal before timing is reported.

\begin{longtable}[]{@{}
  >{\raggedleft\arraybackslash}p{(\columnwidth - 8\tabcolsep) * \real{0.2000}}
  >{\raggedleft\arraybackslash}p{(\columnwidth - 8\tabcolsep) * \real{0.2000}}
  >{\raggedleft\arraybackslash}p{(\columnwidth - 8\tabcolsep) * \real{0.2000}}
  >{\raggedleft\arraybackslash}p{(\columnwidth - 8\tabcolsep) * \real{0.2000}}
  >{\raggedleft\arraybackslash}p{(\columnwidth - 8\tabcolsep) * \real{0.2000}}@{}}
\toprule\noalign{}
\begin{minipage}[b]{\linewidth}\raggedleft
\(\lvert S\rvert\)
\end{minipage} & \begin{minipage}[b]{\linewidth}\raggedleft
parts in \(K(S \cup D)\)
\end{minipage} & \begin{minipage}[b]{\linewidth}\raggedleft
batch reclosure (s)
\end{minipage} & \begin{minipage}[b]{\linewidth}\raggedleft
incremental (s)
\end{minipage} & \begin{minipage}[b]{\linewidth}\raggedleft
speedup
\end{minipage} \\
\midrule\noalign{}
\endhead
\bottomrule\noalign{}
\endlastfoot
250 & 3,546 & 0.035 & 0.006 & 5.9× \\
500 & 6,289 & 0.080 & 0.009 & 9.3× \\
1,000 & 11,496 & 0.154 & 0.015 & 10.2× \\
2,000 & 22,096 & 0.310 & 0.032 & 9.6× \\
4,000 & 43,296 & 0.606 & 0.067 & 7.5--9.6× \\
8,000 & 85,696 & 1.241 & 0.166 & 7.5× \\
\end{longtable}

\begin{figure}
\centering
\includegraphics{fig_benchmark.png}
\caption{Median wall time for integrating a fixed 20-atom delta: batch
reclosure grows linearly with the base; frontier-local incremental
closure stays an order of magnitude below it.}
\end{figure}

Batch cost grows linearly with \(\lvert S\rvert\), as it must --- it
re-derives the whole field. The incremental curve's residual growth is
an artifact of the deliberately unoptimized reference implementation,
which rebuilds its atom index over all atoms once per frontier round;
with premise indexing that term vanishes and the cost is proportional to
the delta closure alone, which is exactly the cost statement
accompanying Theorem 4. The honest headline is the measured one: a
\(6\)--\(10\times\) advantage with no indexing at all, and bit-identical
fingerprints throughout.

\hypertarget{descriptor-cohorts-on-a-natural-corpus-does-proposition-10-bite}{%
\subsubsection{8.2 Descriptor cohorts on a natural corpus --- does
Proposition 10
bite?}\label{descriptor-cohorts-on-a-natural-corpus-does-proposition-10-bite}}

A reproducible natural-language stress proxy: atoms are the 8,631
sentences of \emph{Moby-Dick} (Project Gutenberg \#2701, header/footer
stripped, sentences of ≥ 20 characters); the descriptor of a sentence is
its most frequent non-stopword token of length ≥ 3, ties broken
alphabetically. The map is deliberately crude --- the point is the
\emph{shape} of naturally occurring cohorts, not linguistic quality.

The result: 3,078 descriptors; 637 bindable cohorts (\(n \ge 3\))
covering 5,671 atoms (65.7\%); a heavy Zipf-like tail with maximal
cohort \(n = 268\) (the descriptor \texttt{ahab}). Under the complete
bind semantics the implied block count is \(\sim\!10^{80.7}\) ---
dominated by the maximal cohort alone, and beyond physical storage. The
quotient view of §6.1 materializes 637 blocks; the cohort-object
representation of §6.2 stores 3,078 records.

\begin{figure}
\centering
\includegraphics{fig_cohorts.png}
\caption{Cohort sizes by descriptor rank (log--log): a heavy tail with
one 268-atom cohort. Everything above the dashed \(m_0 = 3\) line is
bindable.}
\end{figure}

The reading is sharp: under natural heavy-tailed descriptor maps, the
complete semantics is not merely expensive but unrealizable in storage,
so the §6 resolutions are forced, not optional. Crucially,
\emph{complete identity} survives: every one of the \(\sim\!10^{80.7}\)
bindings has a well-defined address computable on demand (H4 is a
function, not a table), so the field remains canonical while only a
linear-size representation is ever materialized --- the design principle
of §6 made quantitative.

\hypertarget{verification-status}{%
\subsubsection{8.3 Verification status}\label{verification-status}}

\begin{longtable}[]{@{}
  >{\raggedright\arraybackslash}p{(\columnwidth - 6\tabcolsep) * \real{0.2500}}
  >{\raggedright\arraybackslash}p{(\columnwidth - 6\tabcolsep) * \real{0.2500}}
  >{\raggedright\arraybackslash}p{(\columnwidth - 6\tabcolsep) * \real{0.2500}}
  >{\raggedright\arraybackslash}p{(\columnwidth - 6\tabcolsep) * \real{0.2500}}@{}}
\toprule\noalign{}
\begin{minipage}[b]{\linewidth}\raggedright
Rung
\end{minipage} & \begin{minipage}[b]{\linewidth}\raggedright
Artifact
\end{minipage} & \begin{minipage}[b]{\linewidth}\raggedright
Establishes
\end{minipage} & \begin{minipage}[b]{\linewidth}\raggedright
Result
\end{minipage} \\
\midrule\noalign{}
\endhead
\bottomrule\noalign{}
\endlastfoot
0 & this paper & deductive proofs of all theorems & author-verified \\
1a & \texttt{growth\_check.py} & 15 fixed instance checks with real
SHA-256 fingerprints & 15/15 PASS \\
1b & \texttt{test\_growth\_properties.py} & 14 universally quantified
properties × 60 randomized adversarial instances each, with shrinking &
14/14 PASS \\
2 & \texttt{GrowthClosure.lean} & kernel-checked proofs of Thms 1--3, 5,
6, the Galois adjunction, the descent identity, and Thm 11; zero
imports, zero \texttt{sorry} & compiles clean (Lean 4.30.0) \\
\end{longtable}

The axiom audit, verbatim from the compiler: \texttt{conservativity},
\texttt{galois}, and \texttt{descent} depend on \textbf{no axioms} (pure
constructive logic); \texttt{K\_idem}, \texttt{incremental},
\texttt{ingest\_eq}, \texttt{history\_independent}, and
\texttt{continuity} use only Lean's standard kernel axioms
\texttt{propext} and \texttt{Quot.sound} (via \texttt{funext}), solely
because they are stated as set equalities;
\texttt{Impossibility.impossibility} uses \texttt{propext} only. No
\texttt{Classical.choice} appears anywhere --- the development is
constructive. Deliberately outside the machine-checked perimeter, and
stated as such: Theorem 4 (needs an inductive stage structure; covered
by Rung 1), the discharge of \texttt{Mono\ Φ} from H2 (taken as the
standing hypothesis), and the cryptographic layer of H4 (collision
resistance is an assumption by design).

\begin{center}\rule{0.5\linewidth}{0.5pt}\end{center}

\hypertarget{related-work}{%
\subsection{9. Related work}\label{related-work}}

\textbf{Fixpoints and least models.} Existence and constructivity are
Knaster--Tarski (Tarski 1955) and Kleene; the closure \(K(S)\) is the
least-Herbrand-model construction of logic programming (van Emden \&
Kowalski 1976) transported to typed graphs.

\textbf{Semi-naive evaluation and incremental view maintenance.} Theorem
4 is the semi-naive principle of Datalog evaluation (Bancilhon 1986;
Ullman 1989) restated for content-addressed parts. The contrast with
deletion-capable maintenance (DRed; Gupta, Mumick \& Subrahmanian 1993)
is structural rather than algorithmic: conservativity (Thm. 2) removes
deletions from the model, so the over-delete/rederive machinery is not
optimized away but \emph{absent}.

\textbf{Monotonicity boundaries.} The CALM conjecture (Hellerstein
2010), proved via relational transducer networks (Ameloot, Neven \& Van
den Bussche 2013) and surveyed in Hellerstein \& Alvaro (2020), equates
coordination-free consistency with logical monotonicity; weakened
monotonicities enlarge the coordination-free class (Ameloot, Ketsman,
Neven \& Zinn 2015; Zinn, Green \& Ludäscher 2012). Non-monotonicity of
aggregates in recursion is classical (Ross \& Sagiv 1992/1997). Theorem
11 instantiates this boundary in the content-addressed MDL format; the
positioning paragraph of §5 gives the full dictionary.

\textbf{The chase.} The binding rules are tuple-generating dependencies
firing on every satisfying homomorphism; the closure is a chase (Fagin,
Kolaitis, Miller \& Popa 2003/2005). Content addressing acts as
collision-resistant Skolemization --- the generated null's name is a
function of the generating tuple --- so H4 places the system in the
Skolem/semi-oblivious chase regime (Marnette 2009), where confluence
(Lemma D, Thm. 5) is exactly chase determinism for functional Skolem
terms.

\textbf{MDL.} The admissibility law is two-part-code minimum description
length (Rissanen 1978; Grünwald 2007). Corollary 12 formalizes a known
informal tension: context-sensitive codebooks are globally optimal but
non-monotone; payload-determined codes are locally greedy but
compositional.

\textbf{Content-addressed storage and replication.} H4 and Definition 6
are the Merkle-DAG discipline (Merkle 1987) underlying Git and IPFS;
Theorem 2(iii) is the Merkle diff. Theorems 2--6 jointly state that the
field is a grow-only replicated set (G-Set CRDT) over derivation records
--- union-merge, hence convergent --- which is why ingestion order can
never matter.

\textbf{Merkle-replicated state.} Merkle-CRDTs (Sanjuán, Pöyhtäri,
Teixeira \& Psaras 2020) combine Merkle-DAGs with join-semilattice merge
for convergent replication; the present work differs in studying the
\emph{generation} of the DAG by a rule closure, its growth theory, and
the parsimony impossibility, not replication.

\textbf{Selection as set cover.} Phase 2 of §6.3 is weighted set cover
with the classical greedy guarantee (Chvátal 1979).

\begin{center}\rule{0.5\linewidth}{0.5pt}\end{center}

\hypertarget{limitations}{%
\subsection{10. Limitations}\label{limitations}}

\begin{enumerate}
\def\labelenumi{\arabic{enumi}.}
\tightlist
\item
  \textbf{H2 excludes context-sensitive MDL.} Payload-determined codes
  cannot price a part against the \emph{global} codebook; Theorem 11
  shows this is the exact boundary of the monotone regime, not an
  artifact of the proof technique. Globally optimal compression belongs
  in Phase-2 views (§6.3).
\item
  \textbf{Finiteness requires bounded grade.} \(K(S)\) is finite for
  finite \(S\) only because ascent is strictly grade-increasing over a
  finite tower and admissibility gates the lateral/binding
  combinatorics; unbounded towers reintroduce infinite closures.
\item
  \textbf{Complete semantics is exponential in cohort size.} Proposition
  10 is a theorem about the chosen semantics, and the resolutions of §6
  are mitigations of representation, not of the count itself.
\item
  \textbf{Status.} Working paper; the proofs are elementary but have not
  been independently refereed. The implementation checks the
  quantitative claims on the demonstration tower; the theorems
  themselves are stated for any system satisfying H1--H4.
\end{enumerate}

\begin{center}\rule{0.5\linewidth}{0.5pt}\end{center}

\hypertarget{conclusion}{%
\subsection{11. Conclusion}\label{conclusion}}

A content-addressed rule closure can be grown the way a ledger is grown:
append-only, order-blind, and auditable. Four hypotheses --- finitary
rules, payload-determined admissibility, stratified acyclic generation,
content addressing --- buy four guarantees: nothing is ever retracted
and growth is a Merkle diff; deltas close incrementally from the
frontier with no rederivation; the checkpoint fingerprint is a function
of content, never of history; and streaming is exact. The dynamics those
guarantees induce are themselves lawful: admissibility only switches on,
repair is absorption by seed growth, well-formedness is a cheap
re-checkable predicate rather than an invariant, and the topology only
accumulates cycles. The one thing such a system cannot do is \emph{also}
store exactly the maximal, non-redundant structure: maximality is a
property of the ambient seed, provably inexpressible in
payload-determined side conditions, and insisting on it in storage
forfeits monotonicity, stable identity, and history independence at
once. The trilemma resolves by separation of concerns --- identity and
storage monotone and complete; parsimony in views --- and every claim in
this account is carried by a three-rung ladder: a 666-line
standard-library ledger (\texttt{15/15\ PASS}), a randomized property
suite (14/14 on \textasciitilde840 instances), and a kernel-checked,
constructive Lean 4 development in which the closure algebra and the
impossibility witness are theorems --- with the boundary itself now
positioned where it belongs, as the content-addressed face of CALM.

\hypertarget{references}{%
\subsection{References}\label{references}}

Modern entries were existence-verified (venue, year, pages where shown)
during the v2 prior-art sweep; classical entries are standard.

\begin{enumerate}
\def\labelenumi{\arabic{enumi}.}
\tightlist
\item
  A. Tarski. A lattice-theoretical fixpoint theorem and its
  applications. \emph{Pacific J. Math.} 5(2):285--309, 1955.
\item
  S. C. Kleene. \emph{Introduction to Metamathematics.} North-Holland,
  1952.
\item
  M. H. van Emden, R. A. Kowalski. The semantics of predicate logic as a
  programming language. \emph{J. ACM} 23(4):733--742, 1976.
\item
  F. Bancilhon, R. Ramakrishnan. An amateur's introduction to recursive
  query processing strategies. \emph{SIGMOD} 1986.
\item
  J. D. Ullman. \emph{Principles of Database and Knowledge-Base
  Systems}, Vol. II. Computer Science Press, 1989.
\item
  A. Gupta, I. S. Mumick, V. S. Subrahmanian. Maintaining views
  incrementally. \emph{SIGMOD} 1993, 157--166.
\item
  R. Fagin, P. G. Kolaitis, R. J. Miller, L. Popa. Data exchange:
  semantics and query answering. \emph{ICDT} 2003, 207--224;
  \emph{Theor. Comput. Sci.} 336(1):89--124, 2005.
\item
  B. Marnette. Generalized schema-mappings: from termination to
  tractability. \emph{PODS} 2009, 13--22.
\item
  J. Rissanen. Modeling by shortest data description. \emph{Automatica}
  14(5):465--471, 1978.
\item
  P. D. Grünwald. \emph{The Minimum Description Length Principle.} MIT
  Press, 2007.
\item
  R. C. Merkle. A digital signature based on a conventional encryption
  function. \emph{CRYPTO} 1987.
\item
  V. Chvátal. A greedy heuristic for the set-covering problem.
  \emph{Math. Oper. Res.} 4(3):233--235, 1979.
\item
  M. Shapiro, N. Preguiça, C. Baquero, M. Zawirski. Conflict-free
  replicated data types. \emph{SSS} 2011.
\item
  J. M. Hellerstein. The declarative imperative: experiences and
  conjectures in distributed logic. \emph{SIGMOD Record} 39(1), 2010
  (PODS 2010 keynote).
\item
  T. J. Ameloot, F. Neven, J. Van den Bussche. Relational transducers
  for declarative networking. \emph{PODS} 2011, 283--292; \emph{J. ACM}
  60(2), Art. 15, 2013.
\item
  J. M. Hellerstein, P. Alvaro. Keeping CALM: when distributed
  consistency is easy. \emph{Commun. ACM} 63(9), 2020.
\item
  T. J. Ameloot, B. Ketsman, F. Neven, D. Zinn. Weaker forms of
  monotonicity for declarative networking: a more fine-grained answer to
  the CALM-conjecture. \emph{PODS} 2014, 64--75; \emph{ACM TODS} 40(4),
  Art. 21, 2015.
\item
  D. Zinn, T. J. Green, B. Ludäscher. Win-move is coordination-free
  (sometimes). \emph{ICDT} 2012.
\item
  K. A. Ross, Y. Sagiv. Monotonic aggregation in deductive databases.
  \emph{PODS} 1992, 114--126; \emph{J. Comput. Syst. Sci.} 54(1):79--97,
  1997.
\item
  H. Sanjuán, S. Pöyhtäri, P. Teixeira, Y. Psaras. Merkle-CRDTs:
  Merkle-DAGs meet CRDTs. arXiv:2004.00107, 2020.
\end{enumerate}

\hypertarget{revision-note-v1-v2}{%
\subsection{Revision note --- v1 → v2}\label{revision-note-v1-v2}}

\begin{enumerate}
\def\labelenumi{\arabic{enumi}.}
\tightlist
\item
  \textbf{Theorem 11 repositioned.} Following a prior-art sweep, the
  impossibility result is now stated as a content-addressed,
  MDL-flavored instance of the CALM boundary (positioning paragraph, §5;
  dictionary in §9), with the trilemma formulation, the MDL
  localization, and the mechanized witness as the claimed contributions.
  Abstract and contributions adjusted accordingly.
\item
  \textbf{Verification ladder added} (§8.3): randomized property suite
  (14/14, \textasciitilde840 instances) and a self-contained,
  constructive, \texttt{sorry}-free Lean 4 formalization with verbatim
  axiom audit. v1 predated both.
\item
  \textbf{Empirical scale added} (§8.1--8.2): measured 6--10×
  incremental advantage on bases up to 8,000 atoms with
  fingerprint-equality assertions, and the Moby-Dick cohort experiment
  showing the complete semantics implies \(\sim\!10^{80.7}\) bindings on
  natural heavy-tailed descriptors --- making §6 quantitatively forced.
\item
  \textbf{Bibliography formalized.} Informal author-year mentions
  replaced by a 20-entry reference list; eleven modern entries
  existence-verified during the sweep.
\item
  \textbf{Related work extended} with the CALM cluster, weakened
  monotonicity, Ross--Sagiv aggregation, Marnette's chase analysis, and
  Merkle-CRDTs.
\end{enumerate}

\begin{center}\rule{0.5\linewidth}{0.5pt}\end{center}

\emph{Artifacts:} \texttt{growth\_check.py} (ledger) ·
\texttt{test\_growth\_properties.py} (property suite) ·
\texttt{GrowthClosure.lean} (kernel-checked proofs) ·
\texttt{bench\_growth.py}, \texttt{corpus\_cohorts.py} (§8 measurements)
· all deterministic and reproducible from the repository.

\end{document}
